\centerline{\bf \Huge Вступительная проверочная работа}

\begin{flushright}
        \scriptsize{\textbf{Именно математика даёт надёжнейшие правила: тому кто им следует — тому не опасен обман чувств.\\
         Леонард Эйлер}}
        
\end{flushright}

\problem{\textbf{1}} Решите систему из двух линейных алгебраических уравнений(СЛАУ):
\begin{equation*}
 \begin{cases}
   100y - 99x = -777
   \\
   297x - 300y = 2331
 \end{cases}
\end{equation*}

\begin{flushright} \textbf{3 балла}\\ 
\end{flushright}

\problem{\textbf{2}} Решите систему из трёх линейных алгебраических уравнений(СЛАУ):
\begin{equation*}
 \begin{cases}
   x + 2y + 3z = 12
   \\
   2x - y + z = -1
   \\
   3x + 7y - 2z = 5
   
 \end{cases}
\end{equation*}


\problem{\textbf{3}} Чему равно следующее выражение: 

$50\cdot\Bigr(\dfrac{(1,238+2,762) \cdot 0,1}{(36,487-34,237): 2,8125}+\dfrac{(4,36-1,16) \cdot 0,3125}{0,2 \cdot(47,8-45,55): 0,225}\Bigl)\cdot$ 

$\cdot \Bigr(39,86-\dfrac{8 \dfrac{4}{5}}{6,934+38 \dfrac{13}{40}: 2 \dfrac{1}{2} \cdot 0,2}-17\Bigl)$ = ?

\begin{flushright} \textbf{6 баллов}\\ 
\end{flushright}


\problem{\textbf{4}} Прямая $Q$ проходит через точки $A(11\dfrac{13}{4};-14,5)$ и $B(-9,25;2\dfrac{101}{4})$.

Найдите уравнение, прямой проходящей через точку $C(-11,4;2,5)$, перпендикулярной прямой $Q$.


\problem{\textbf{5}} Назовём трёхзначное число интересным, если хотя бы одна его
цифра делится на 3. Какое наибольшее количество подряд идущих интересных чисел может
быть? Приведите пример таких чисел и докажите, что больше быть не может!

\begin{flushright} \textbf{9 баллов}\\ 
\end{flushright}


\problem{\textbf{8.1}} Постройте (по точкам) график функции кривой второго порядка $y=x^2-10x+25+4$.
\begin{flushright} \textbf{6 баллов}\\ 
\end{flushright}

\newpage

\hspace{1mm}

\problem{\textbf{8.2}} \textbf{\textit{Новые множества чисел:}}

Рациональное число ($\mathbb{Q}$ - множество рациональных чисел, с французского 
$qutient$ - отношение) - это такое число, которое \textbf{МОЖНО} представить в виде несократимой обыкновенной дроби $\dfrac{m}{n}$ , где $m\in \mathbb{Z}, n\in \mathbb{N}$. ($\mathbb{Z}$ - множество целых чисел, $\mathbb{N}$ - множество натуральных чисел)

Иррациональное число ($\mathbb{I}$ - множество иррациональных чисел, с латинского 

$irrationalis$ - не выражаемый через отношение) - это такое число, которое \textbf{НЕЛЬЗЯ} представить в виде несократимой обыкновенной дроби $\dfrac{m}{n}$ , где $m\in \mathbb{Z}, n\in \mathbb{N}$.

Определите, являются рациональными или иррациональными следующие числа:

а) 12,3456789?

б) $X$? - такое, что выполняется равенство $X^2=7$

в) 2,(5)?

\problem{\textbf{8.3}} Решением неравенства $\dfrac{\left(x-5\right)\left(x+6\right)}{\left(x+8\right)\left(x+5\right)}\ge0$ 

Является следующее множество чисел $x\in(-\infty;-8) \cup[-6;-5)\cup[5;+\infty)$.

Определите множество чисел, являющихся решением следующего неравенства:

$\dfrac{\left(x+1\right)\left(x-9\right)\left(x+11\right)}{\left(x-2\right)\left(x+7\right)}\le0$

\problem{\textbf{8.4}} \textbf{\textit{Теорема Виета:}}

Пусть $x_1,x_2$ корни квадратного уравнения $ax^2+bx+c=0;$

Тогда верна следующая система:
\begin{equation*}
 \begin{cases}
   x_1+x_2=-\dfrac{b}{a}
   \\
   x_1\cdot x_2=\dfrac{c}{a}
 \end{cases}
\end{equation*}

Используя теорему Виета, найдите корни (они $\in \mathbb{Z}$) следующего квадратного уравнения: $2x^{2}-30x+112=0$
\begin{flushright} \textbf{15 баллов}\\ 
\end{flushright}


